%%%%%%%%%%%%%%%%%%%%%%%%%%%%%%%%%%%%%%%%%%%%%%%%%%%%%%%%%%%%%
% BibLaTeX 설정
% - style=apa: APA 7판 기준 (biblatex-apa 패키지)
% - 국문/영문 분리 출력:
%     \printbibliography[check=korean, title={참고문헌}]
%     \printbibliography[check=notkorean, heading=none]
%%%%%%%%%%%%%%%%%%%%%%%%%%%%%%%%%%%%%%%%%%%%%%%%%%%%%%%%%%%%%

\usepackage[
  backend=biber,
  style=apa,
  sorting=nyt,
  dashed=false,
  maxbibnames=10,
  minbibnames=6,
  maxcitenames=3,
  mincitenames=1
]{biblatex}
\addbibresource{refs.bib}

%------------------------------------------------------------
% 공통 포맷 설정
%------------------------------------------------------------

% "in:" 제거
\renewbibmacro{in:}{}

% 페이지 범위 앞의 'pp.' 제거
\DeclareFieldFormat{pages}{#1}

% article, inproceedings 제목: 따옴표 없이 그대로 (APA 기본값 유지)
%\DeclareFieldFormat[article,inproceedings]{title}{#1}

% 공저자 구분자: & 앞뒤 공백 (영문 기본값)
\renewcommand{\finalnamedelim}{\ \&\ }

%------------------------------------------------------------
% 국문 참고문헌 전용 설정 (langid=korean 기준)
%------------------------------------------------------------

% 학위논문 유형 한국어 표기: preamble에서 한 번만 선언
\DefineBibliographyStrings{english}{%
  phdthesis = {박사학위논문},
  mathesis  = {석사학위논문},
}

%\AtEveryBibitem{%
%  \iffieldequalstr{langid}{korean}{%
%    % 국문: 저자 구분자 쉼표로 변경
%    \renewcommand{\multinamedelim}{,\space}%
%    \renewcommand{\finalnamedelim}{,\space}%
%  }{}%
%}
%\AtEveryBibitem{%
%  \iffieldequalstr{langid}{korean}{%
%    \renewcommand{\multinamedelim}{,\space}%
%    \renewcommand{\finalnamedelim}{,\space}%
%    \renewcommand{\mkbibemph}[1]{##1}%   % 이탤릭 무력화만 유지
%    \renewcommand{\emph}[1]{##1}%
%  }{}%
%}
\AtEveryBibitem{%
  \iffieldequalstr{langid}{korean}{%
    \renewcommand{\multinamedelim}{,\space}%
%    \renewcommand{\finalnamedelim}{,\space}%
    % biblatex-apa의 & 삽입 매크로 무력화
    \renewcommand{\andothersdelim}{\addcomma\space}%
    \DeclareDelimFormat{andothers}{\addcomma\space}%
    \renewbibmacro*{name:andothers}{%
      \ifboolexpr{
        test {\ifnumequal{\value{listcount}}{\value{liststop}}}
        and
        test \ifmorenames
      }
        {\printdelim{andothersdelim}\bibstring{andothers}}
        {}}%
  }{}%
}
\usepackage{xpatch}

% biblatex-apa의 finalnamedelim을 langid에 따라 분기
\DeclareDelimFormat{finalnamedelim}{%
  \iffieldequalstr{langid}{korean}%
    {\addcomma\space}%
    {\addspace\&\space}%
}
%\AtEveryBibitem{%
%  \iffieldequalstr{langid}{korean}{%
%    \renewcommand{\multinamedelim}{,\space}%
%    \renewcommand{\finalnamedelim}{,\space}%
%    % 이탤릭 → 볼드 대체
%    \DeclareFieldFormat[article]{title}{#1}%               % 제목: 일반체
%    \DeclareFieldFormat[article]{journaltitle}{\textbf{#1}}% 학술지명: 볼드
%    \DeclareFieldFormat[article]{volume}{\textbf{#1}}%      % 권: 볼드
%    \DeclareFieldFormat[article]{number}{\mkbibparens{#1}}% % 호: 괄호, 일반체
%  }{}%
%}
% \DeclareFieldFormat는 \AtEveryBibitem 밖, preamble에서 선언
\DeclareFieldFormat[article]{journaltitle}{%
  \iffieldequalstr{langid}{korean}{\textbf{#1}}{\mkbibemph{#1}}%
}
\DeclareFieldFormat[article]{volume}{%
  \iffieldequalstr{langid}{korean}{\textbf{#1}}{\mkbibemph{#1}}%
}
\DeclareFieldFormat[article]{number}{%
  \mkbibparens{#1}%
}
% preamble 레벨 (AtEveryBibitem 밖)
\DeclareFieldFormat[article]{title}{%
  \iffieldequalstr{langid}{korean}{#1}{\mkbibquote{#1}}%
}
%------------------------------------------------------------
% \textcite 사용 시 저자명과 연도 사이 공백 제거
%------------------------------------------------------------

\DeclareDelimFormat[textcite]{nameyeardelim}{}

%------------------------------------------------------------
% 국문/영문 분리 필터
%------------------------------------------------------------

\defbibcheck{korean}{%
  \iffieldequalstr{langid}{korean}{}{\skipentry}%
}
\defbibcheck{notkorean}{%
  \iffieldequalstr{langid}{korean}{\skipentry}{}%
}
